%! Unit 1 — Functions and Change — Unit Cover Page (shared body)
%! Inputted by BOTH unit_cover/main.tex and unit_cover_key/main.tex so that
%! page 1 of the student packet and page 1 of the key packet cannot drift.
%! Edit the cover HERE, never in a wrapper.
\thispagestyle{empty}

% Full-bleed plum banner
\begin{tikzpicture}[remember picture, overlay]
  \fill[plum] (current page.north west)
    rectangle ([yshift=-1.15in]current page.north east);
\end{tikzpicture}
\vspace{-0.82in}

\begin{center}
  {\color{white}\textbf{\LARGE Precalculus}}\\[6pt]
  {\color{lilacmid}\Large\textbf{Unit 1: Functions and Change}}\\[4pt]
  {\color{white}\small 8 Lessons + Unit Test \quad|\quad Class Periods: 18--22}
\end{center}

\vspace{0.20in}

\begin{tcolorbox}[colback=lilac, colframe=plum, boxrule=0.8pt, arc=2mm,
    left=3mm, right=3mm, top=2mm, bottom=2mm]
\small\textbf{Unit Overview.}
This unit is about two quantities that change together, and about the single
idea that lets us talk about them: a \textbf{function} is a rule that pairs each
input with exactly one output. The first half of the unit learns to
\emph{describe} such a rule --- reading it in notation, in a table, and in a
graph; naming where it increases and decreases; deciding whether it bends up or
down; and measuring how fast it changes with an average rate of change. The
second half learns to \emph{build} one: shifting and stretching a parent
function, feeding one function into another, and running a function backward to
recover the input that produced a given output. The unit closes by putting both
halves to work --- choosing a model that fits a situation, constructing it from
a constraint, and stating the domain and the assumptions it rests on. Everything
in the seven units that follow is a particular family of functions; this unit is
the vocabulary they will all be discussed in.
\end{tcolorbox}

\vspace{0.13in}

\begin{skillbox}[Lessons in This Unit]{goldbox}
\small
\renewcommand{\arraystretch}{1.42}
\begin{tabularx}{\linewidth}{@{} c >{\bfseries\raggedright\arraybackslash}p{0.40\linewidth} X @{}}
\rowcolor{lilac}
\textbf{\#} & \textbf{Lesson Title} & \textbf{Focus} \\
\midrule
1.0 & Introduction to Functions and Change & Inputs, outputs, and what makes a rule a function \\
1.1 & Function Fundamentals & Notation, evaluating vs.\ solving, domain and range \\
1.2 & Change in Tandem & Increasing, decreasing, concavity, zeros \\
1.3 & Rates of Change & Average rate of change, units, first and second differences \\
1.4 & Transformations of Functions & Changing the question vs.\ changing the answer \\
1.5 & Composition of Functions & The handoff: $f(g(x))$, order, decomposition \\
1.6 & Inverse Functions & Giving back the question; one-to-one; swap and solve \\
1.7 & Function Model Selection and Construction & Picking the question, building and using a model \\
\midrule
\rowcolor{lilac}
    & \textbf{Sample Test} & A study copy of the unit test, with an answer key \\
\end{tabularx}
\end{skillbox}

\vspace{0.13in}

\begin{skillbox}[Mathematical Practices --- Unit 1 Focus]{greenbox}
\small
\renewcommand{\arraystretch}{1.32}
\begin{tabularx}{\linewidth}{@{} >{\leavevmode\bfseries\color{plum}}p{0.11\linewidth} X @{}}
MP 1.C & Construct new functions using transformations, compositions, inverses, or regressions. \\
MP 2.A & Identify information from graphical, numerical, analytical, and verbal representations. \\
MP 2.B & Construct equivalent graphical, numerical, analytical, and verbal representations. \\
MP 3.A & Describe the characteristics of a function with varying levels of precision. \\
MP 3.C & Support conclusions or choices with a logical rationale or appropriate data. \\
\end{tabularx}
\end{skillbox}

\vspace{0.13in}

\begin{remindbox}
\small
\textbf{Two questions carry the whole unit:} \emph{How do we describe the way one
quantity changes with respect to another?} and \emph{How do we build a new
function out of ones we already have?} Keep them in view --- every lesson here
answers one of the two, and the unit test asks you to do both.
\end{remindbox}
